\section{Experimental Setup}
\label{sec:setup}

% DRAFT (Part T, 2026-11-27). Sources: docs/dataset_notes.md,
% docs/labels_notes.md, docs/geometry_notes.md, docs/rerender_notes.md,
% docs/MORNING_REPORT.md, docs/experiment_registry.md, policy/*.py.

\draftnote{DRAFT. Protocol as built. No trained-policy success number exists
yet; see ``What has not been run''.}

\paragraph{Simulator and tasks.} We use RoboCasa v1.0.1
\citep{nasiriany2024robocasa} at commit \texttt{425684d} with robosuite
\texttt{5ce6643} and MuJoCo 3.3.1, on the target split of RoboCasa365
\citep{nasiriany2026robocasa365}. The robot is a Panda arm on a mobile base
(\texttt{PandaOmron}), controlled at 20\,Hz with delta end-effector commands
in the base frame. We use the 18 atomic-seen tasks that \dem{} reports on
\citep{sun2026dem}: five pick-and-place tasks, \texttt{CoffeeSetupMug},
doors and lids, a drawer, a dishwasher rack, buttons, knobs, levers and
\texttt{NavigateKitchen}. Three of them serve as fast screening tasks:
\texttt{PickPlaceCounterToCabinet}, \texttt{OpenDrawer} and
\texttt{TurnOnMicrowave} (one pick-and-place, one articulated fixture, one
button). Kitchens come from the target layouts and styles 1 to 10, and
objects from the target object split.

\paragraph{Data.} We train on the human demonstrations of the target split
in LeRobot format: 500 to 543 per task, 9,126 episodes and 2.23 million
frames in total. Each episode has left, right and wrist videos at
$256\times256$, the 16-number state, the 12-number action, the instruction,
and saved simulator states. We use the left and right agent views.
\draftnote{The DEM README's example uses a scene camera plus the wrist; a
left-plus-wrist run is training to decide this.}
Hindsight labels (\cref{sec:labels}) exist for every episode of all 18 tasks.
Five pick-and-place tasks and \texttt{CoffeeSetupMug} are labeled from
gripper close and open in every episode. Pushing and pressing tasks fall back
in part or in full to the success frame (for example, all 502 episodes of
\texttt{OpenStandMixerHead}, 51 of 543 of \texttt{TurnOnMicrowave}).
Checked in the simulator on 20 demonstrations per task, the grasp point lies
within the object's extent in all 120 pick-style demonstrations, with a
median distance to the object's origin of 1.6 to 5.6\,cm by task.

\paragraph{Normalization.} We follow the \dem{} research code in two stages.
Stage 1 maps each task's 1st to 99th percentile of each state and action
dimension to $[-1,1]$ (constant dimensions map to $-1$ and are restored
exactly). Stage 2 z-scores with statistics over all training tasks, inside
the policy. Both are saved with each checkpoint.

\paragraph{Training.} The recipe follows the \dem{} release: AdamW with
learning rate $10^{-4}$ for the head and $5\times10^{-5}$ for the vision
encoder, weight decay $10^{-4}$, $\beta=(0.9, 0.95)$, a cosine schedule,
gradient clipping at 1.0, batch size 64, bf16 autocast for the vision encoder
and fp32 for the head and loss. NeoBERT is frozen and its tokens are cached.
Screening runs train for 20k steps on the three smoke tasks; the full run
trains for 300k steps on all 18 tasks, as in \dem{}. Every format uses the
same steps, batch, data order and seed (seed 0); two more seeds (1 and 2)
measure run-to-run noise for stock \dem{}.

\paragraph{Evaluation.} Each episode resets the environment with seed
$1000+i$ for episode $i$, so all formats and conditions see the same kitchens
and objects. The policy predicts 16 actions and executes 8, then predicts
again. Success is checked after every step with the simulator's own check,
and the episode stops at success or at the task's horizon from RoboCasa's
registry (450 to 1,050 steps, 22.5 to 52.5\,s of robot time). Screening
checkpoints are scored every 5k steps with 50 episodes per task; anything
that may go in the paper is scored with 100 episodes per task. We report
success rates with 95\% Wilson score intervals and mark comparisons whose
intervals overlap. Every run writes one CSV row per episode, with the code
commit.

\paragraph{Checks before any number counts.} Recorded demonstration actions,
replayed from the saved start state, succeed in 30 of 30 demonstrations on
the three smoke tasks. The same actions run through our evaluation loop
succeed in 30 of 30; the state at reset matches the recording to $10^{-8}$,
and rendered images differ from the video frames by 1.7 to 2.5 grey levels
(out of 255) on average. An untrained head scores 0 of 28 (95\% interval 0 to
12\%).

\paragraph{Perfect commands.} To score a format without an agent, we build
the command from simulator ground truth at reset and treat it as an ideal
agent's command. The text uses the same template and object names as the
training labels. For pick-and-place tasks the grasp point is the object's
body origin and the place point is the destination (cabinet interior centre,
pan, plate, coffee receptacle); for fixture tasks it is the task's part
(handle, knob, button, lever) nearest the gripper. These points do not match
the training labels exactly: labels are where the gripper closed, perfect
commands are body origins, a median gap of 1.6 to 5.6\,cm by task and up to
about 10\,cm. \draftnote{Known mismatch. It biases perfect-command scores of
point formats downward; not yet measured.} We also score controlled errors:
points shifted by a fixed distance, the grasp point on a distractor object (to
test whether the policy follows the point), dropped fields, and renamed
objects.

\paragraph{Held-out objects.} For pick-and-place we hold out object
categories with a fixed seed. A demonstration is used for training only if no
object in its scene, target or distractor, is held out. For
\texttt{PickPlaceCounterToCabinet}, 16 of 106 categories are held out,
leaving 328 training demonstrations, 35 validation and 36 locked test
demonstrations (103 dropped because a held-out object appears as a
distractor). Held-out evaluation replays those demonstrations' initial
scenes. The locked test half is scored only for results going into the
paper. \draftnote{Other pick-and-place tasks: split not yet built.}

\paragraph{Agent.} The commander (\cref{sec:agent}) runs Claude Opus 5.5
through a command-line coding agent on a subscription, and we log turns,
input and output tokens, and wall-clock seconds per episode. Pixels are
lifted with rendered depth; on a single scene a lifted mask-centre pixel lands
3.2 and 3.5\,cm from the object origin in the left and right views.

\paragraph{Hardware.} Training runs on single GPUs of a university cluster.
On an H100 80\,GB at batch 64 with two cameras, training runs at 3.27 steps
per second (83\% GPU busy, no data wait, 32\,GB peak), so 20k steps take
about 1.7 hours and 300k steps about 25.5 hours. An L40S runs at 1.16 steps
per second (20k steps in about 4.8 hours).

\paragraph{What has not been run.} As of this draft: no trained checkpoint
has been scored. Stock \dem{} is training on the three smoke tasks (three
seeds and a left-plus-wrist camera variant) and on all 18 tasks toward 300k
steps. Of the format runs, \texttt{kv\_text} is training on the smoke tasks;
\texttt{points\_3d\_rgb}, \texttt{points\_2d} and \texttt{pose\_path} are
queued. \texttt{points\_3d} with point-map input is not implemented in the
training loader. Depth re-rendering has been checked on 10 demonstrations of
one task. The commander has run one episode, with a 50-step checkpoint, as a
smoke test. The agent-alone baseline has not been run.

\paragraph{Conditions and metrics.} \texttt{policy\_only} (perfect or fixed
commands), \texttt{agent\_only}, and \texttt{agent\_plus\_policy}. We report
success rate, seconds per episode, LLM calls and tokens per episode.
